\section*{Item 1}\label{it:1}

\textbf{Enunciado}: um conversor buck tem uma tensão de entrada que varia entre $10$ e $15$ V e uma corrente de carga que varia entre $2{,}0$ A e $3{,}5$ A. A tensão de saída é $5$ V. Para uma frequência de comutação de $200$ kHz, determine a indutância mínima necessária para fornecer corrente contínua para todas as possibilidades de operação.

\medskip
\textbf{Resolução}: partindo da equação que define o limite entre os modos de operação CCM e DCM para um conversor Buck \cite{Eri:20}
\begin{equation}
    K = \frac{2L}{RT_s} = 1 - D
\end{equation}

Para o CCM, tem-se que $K > 1-D$. Portanto, quanto menor o valor de $D$, maior o limitante inferior de K. O pior caso então é obtido com $V_g = 15$, que implica $D = V/V_g = 0{,}333$. Além disso, quanto maior o valor de $R$, maior a tendência do conversor operar no DCM. Então, a pior condição é $R = 5/2 = 2{,}5$. Finalmente, pode-se definir o menor valor de $L$ por
\begin{equation}\label{eq:ex1_1}
\begin{aligned}
    L > \frac{RT_s}{2}(1-D) &= \frac{2{,}5(2\times10^5)}{2}(1-0{,}3) \\
    &\approx 4{,}1667 \ \mu\mathrm{H}
\end{aligned}
\end{equation}

Este valor foi verificado por meio de simulação utilizando o \textit{software} ngspice. A Figura~\ref{fig:ex1} mostra a tensão de saída e a corrente no indutor, que atinge o valor nulo ao final do período de comutação, conforme previsto pela condição~\eqref{eq:ex1_1}. Para este experimento, foi obtido $\Delta i_L = 2$ A e calculado $C = 2{,}5 \ \mu$F para um \textit{ripple} arbitrariamente escolhido como sendo $0{,}5$ V~\cite{codm_ex1}.
\begin{figure}[htb!]
    \centering
    %% Creator: Matplotlib, PGF backend
%%
%% To include the figure in your LaTeX document, write
%%   \input{<filename>.pgf}
%%
%% Make sure the required packages are loaded in your preamble
%%   \usepackage{pgf}
%%
%% Also ensure that all the required font packages are loaded; for instance,
%% the lmodern package is sometimes necessary when using math font.
%%   \usepackage{lmodern}
%%
%% Figures using additional raster images can only be included by \input if
%% they are in the same directory as the main LaTeX file. For loading figures
%% from other directories you can use the `import` package
%%   \usepackage{import}
%%
%% and then include the figures with
%%   \import{<path to file>}{<filename>.pgf}
%%
%% Matplotlib used the following preamble
%%   \def\mathdefault#1{#1}
%%   \everymath=\expandafter{\the\everymath\displaystyle}
%%   \IfFileExists{scrextend.sty}{
%%     \usepackage[fontsize=8.000000pt]{scrextend}
%%   }{
%%     \renewcommand{\normalsize}{\fontsize{8.000000}{9.600000}\selectfont}
%%     \normalsize
%%   }
%%   
%%           \usepackage{amsmath}
%%           \usepackage{amssymb}
%%           \usepackage{mathtools}
%%           \usepackage{dsfont}
%%           \providecommand{\mathdefault}[1]{#1}
%%       
%%   \makeatletter\@ifpackageloaded{underscore}{}{\usepackage[strings]{underscore}}\makeatother
%%
\begingroup%
\makeatletter%
\begin{pgfpicture}%
\pgfpathrectangle{\pgfpointorigin}{\pgfqpoint{3.334740in}{2.322654in}}%
\pgfusepath{use as bounding box, clip}%
\begin{pgfscope}%
\pgfsetbuttcap%
\pgfsetmiterjoin%
\definecolor{currentfill}{rgb}{1.000000,1.000000,1.000000}%
\pgfsetfillcolor{currentfill}%
\pgfsetlinewidth{0.000000pt}%
\definecolor{currentstroke}{rgb}{1.000000,1.000000,1.000000}%
\pgfsetstrokecolor{currentstroke}%
\pgfsetdash{}{0pt}%
\pgfpathmoveto{\pgfqpoint{0.000000in}{0.000000in}}%
\pgfpathlineto{\pgfqpoint{3.334740in}{0.000000in}}%
\pgfpathlineto{\pgfqpoint{3.334740in}{2.322654in}}%
\pgfpathlineto{\pgfqpoint{0.000000in}{2.322654in}}%
\pgfpathlineto{\pgfqpoint{0.000000in}{0.000000in}}%
\pgfpathclose%
\pgfusepath{fill}%
\end{pgfscope}%
\begin{pgfscope}%
\pgfsetbuttcap%
\pgfsetmiterjoin%
\definecolor{currentfill}{rgb}{1.000000,1.000000,1.000000}%
\pgfsetfillcolor{currentfill}%
\pgfsetlinewidth{0.000000pt}%
\definecolor{currentstroke}{rgb}{0.000000,0.000000,0.000000}%
\pgfsetstrokecolor{currentstroke}%
\pgfsetstrokeopacity{0.000000}%
\pgfsetdash{}{0pt}%
\pgfpathmoveto{\pgfqpoint{0.514740in}{1.415315in}}%
\pgfpathlineto{\pgfqpoint{3.234740in}{1.415315in}}%
\pgfpathlineto{\pgfqpoint{3.234740in}{2.222654in}}%
\pgfpathlineto{\pgfqpoint{0.514740in}{2.222654in}}%
\pgfpathlineto{\pgfqpoint{0.514740in}{1.415315in}}%
\pgfpathclose%
\pgfusepath{fill}%
\end{pgfscope}%
\begin{pgfscope}%
\pgfpathrectangle{\pgfqpoint{0.514740in}{1.415315in}}{\pgfqpoint{2.720000in}{0.807339in}}%
\pgfusepath{clip}%
\pgfsetbuttcap%
\pgfsetroundjoin%
\pgfsetlinewidth{0.501875pt}%
\definecolor{currentstroke}{rgb}{0.690196,0.690196,0.690196}%
\pgfsetstrokecolor{currentstroke}%
\pgfsetstrokeopacity{0.250000}%
\pgfsetdash{{1.850000pt}{0.800000pt}}{0.000000pt}%
\pgfpathmoveto{\pgfqpoint{0.638375in}{1.415315in}}%
\pgfpathlineto{\pgfqpoint{0.638375in}{2.222654in}}%
\pgfusepath{stroke}%
\end{pgfscope}%
\begin{pgfscope}%
\pgfsetbuttcap%
\pgfsetroundjoin%
\definecolor{currentfill}{rgb}{0.000000,0.000000,0.000000}%
\pgfsetfillcolor{currentfill}%
\pgfsetlinewidth{0.803000pt}%
\definecolor{currentstroke}{rgb}{0.000000,0.000000,0.000000}%
\pgfsetstrokecolor{currentstroke}%
\pgfsetdash{}{0pt}%
\pgfsys@defobject{currentmarker}{\pgfqpoint{0.000000in}{-0.048611in}}{\pgfqpoint{0.000000in}{0.000000in}}{%
\pgfpathmoveto{\pgfqpoint{0.000000in}{0.000000in}}%
\pgfpathlineto{\pgfqpoint{0.000000in}{-0.048611in}}%
\pgfusepath{stroke,fill}%
}%
\begin{pgfscope}%
\pgfsys@transformshift{0.638375in}{1.415315in}%
\pgfsys@useobject{currentmarker}{}%
\end{pgfscope}%
\end{pgfscope}%
\begin{pgfscope}%
\pgfpathrectangle{\pgfqpoint{0.514740in}{1.415315in}}{\pgfqpoint{2.720000in}{0.807339in}}%
\pgfusepath{clip}%
\pgfsetbuttcap%
\pgfsetroundjoin%
\pgfsetlinewidth{0.501875pt}%
\definecolor{currentstroke}{rgb}{0.690196,0.690196,0.690196}%
\pgfsetstrokecolor{currentstroke}%
\pgfsetstrokeopacity{0.250000}%
\pgfsetdash{{1.850000pt}{0.800000pt}}{0.000000pt}%
\pgfpathmoveto{\pgfqpoint{1.132921in}{1.415315in}}%
\pgfpathlineto{\pgfqpoint{1.132921in}{2.222654in}}%
\pgfusepath{stroke}%
\end{pgfscope}%
\begin{pgfscope}%
\pgfsetbuttcap%
\pgfsetroundjoin%
\definecolor{currentfill}{rgb}{0.000000,0.000000,0.000000}%
\pgfsetfillcolor{currentfill}%
\pgfsetlinewidth{0.803000pt}%
\definecolor{currentstroke}{rgb}{0.000000,0.000000,0.000000}%
\pgfsetstrokecolor{currentstroke}%
\pgfsetdash{}{0pt}%
\pgfsys@defobject{currentmarker}{\pgfqpoint{0.000000in}{-0.048611in}}{\pgfqpoint{0.000000in}{0.000000in}}{%
\pgfpathmoveto{\pgfqpoint{0.000000in}{0.000000in}}%
\pgfpathlineto{\pgfqpoint{0.000000in}{-0.048611in}}%
\pgfusepath{stroke,fill}%
}%
\begin{pgfscope}%
\pgfsys@transformshift{1.132921in}{1.415315in}%
\pgfsys@useobject{currentmarker}{}%
\end{pgfscope}%
\end{pgfscope}%
\begin{pgfscope}%
\pgfpathrectangle{\pgfqpoint{0.514740in}{1.415315in}}{\pgfqpoint{2.720000in}{0.807339in}}%
\pgfusepath{clip}%
\pgfsetbuttcap%
\pgfsetroundjoin%
\pgfsetlinewidth{0.501875pt}%
\definecolor{currentstroke}{rgb}{0.690196,0.690196,0.690196}%
\pgfsetstrokecolor{currentstroke}%
\pgfsetstrokeopacity{0.250000}%
\pgfsetdash{{1.850000pt}{0.800000pt}}{0.000000pt}%
\pgfpathmoveto{\pgfqpoint{1.627467in}{1.415315in}}%
\pgfpathlineto{\pgfqpoint{1.627467in}{2.222654in}}%
\pgfusepath{stroke}%
\end{pgfscope}%
\begin{pgfscope}%
\pgfsetbuttcap%
\pgfsetroundjoin%
\definecolor{currentfill}{rgb}{0.000000,0.000000,0.000000}%
\pgfsetfillcolor{currentfill}%
\pgfsetlinewidth{0.803000pt}%
\definecolor{currentstroke}{rgb}{0.000000,0.000000,0.000000}%
\pgfsetstrokecolor{currentstroke}%
\pgfsetdash{}{0pt}%
\pgfsys@defobject{currentmarker}{\pgfqpoint{0.000000in}{-0.048611in}}{\pgfqpoint{0.000000in}{0.000000in}}{%
\pgfpathmoveto{\pgfqpoint{0.000000in}{0.000000in}}%
\pgfpathlineto{\pgfqpoint{0.000000in}{-0.048611in}}%
\pgfusepath{stroke,fill}%
}%
\begin{pgfscope}%
\pgfsys@transformshift{1.627467in}{1.415315in}%
\pgfsys@useobject{currentmarker}{}%
\end{pgfscope}%
\end{pgfscope}%
\begin{pgfscope}%
\pgfpathrectangle{\pgfqpoint{0.514740in}{1.415315in}}{\pgfqpoint{2.720000in}{0.807339in}}%
\pgfusepath{clip}%
\pgfsetbuttcap%
\pgfsetroundjoin%
\pgfsetlinewidth{0.501875pt}%
\definecolor{currentstroke}{rgb}{0.690196,0.690196,0.690196}%
\pgfsetstrokecolor{currentstroke}%
\pgfsetstrokeopacity{0.250000}%
\pgfsetdash{{1.850000pt}{0.800000pt}}{0.000000pt}%
\pgfpathmoveto{\pgfqpoint{2.122012in}{1.415315in}}%
\pgfpathlineto{\pgfqpoint{2.122012in}{2.222654in}}%
\pgfusepath{stroke}%
\end{pgfscope}%
\begin{pgfscope}%
\pgfsetbuttcap%
\pgfsetroundjoin%
\definecolor{currentfill}{rgb}{0.000000,0.000000,0.000000}%
\pgfsetfillcolor{currentfill}%
\pgfsetlinewidth{0.803000pt}%
\definecolor{currentstroke}{rgb}{0.000000,0.000000,0.000000}%
\pgfsetstrokecolor{currentstroke}%
\pgfsetdash{}{0pt}%
\pgfsys@defobject{currentmarker}{\pgfqpoint{0.000000in}{-0.048611in}}{\pgfqpoint{0.000000in}{0.000000in}}{%
\pgfpathmoveto{\pgfqpoint{0.000000in}{0.000000in}}%
\pgfpathlineto{\pgfqpoint{0.000000in}{-0.048611in}}%
\pgfusepath{stroke,fill}%
}%
\begin{pgfscope}%
\pgfsys@transformshift{2.122012in}{1.415315in}%
\pgfsys@useobject{currentmarker}{}%
\end{pgfscope}%
\end{pgfscope}%
\begin{pgfscope}%
\pgfpathrectangle{\pgfqpoint{0.514740in}{1.415315in}}{\pgfqpoint{2.720000in}{0.807339in}}%
\pgfusepath{clip}%
\pgfsetbuttcap%
\pgfsetroundjoin%
\pgfsetlinewidth{0.501875pt}%
\definecolor{currentstroke}{rgb}{0.690196,0.690196,0.690196}%
\pgfsetstrokecolor{currentstroke}%
\pgfsetstrokeopacity{0.250000}%
\pgfsetdash{{1.850000pt}{0.800000pt}}{0.000000pt}%
\pgfpathmoveto{\pgfqpoint{2.616558in}{1.415315in}}%
\pgfpathlineto{\pgfqpoint{2.616558in}{2.222654in}}%
\pgfusepath{stroke}%
\end{pgfscope}%
\begin{pgfscope}%
\pgfsetbuttcap%
\pgfsetroundjoin%
\definecolor{currentfill}{rgb}{0.000000,0.000000,0.000000}%
\pgfsetfillcolor{currentfill}%
\pgfsetlinewidth{0.803000pt}%
\definecolor{currentstroke}{rgb}{0.000000,0.000000,0.000000}%
\pgfsetstrokecolor{currentstroke}%
\pgfsetdash{}{0pt}%
\pgfsys@defobject{currentmarker}{\pgfqpoint{0.000000in}{-0.048611in}}{\pgfqpoint{0.000000in}{0.000000in}}{%
\pgfpathmoveto{\pgfqpoint{0.000000in}{0.000000in}}%
\pgfpathlineto{\pgfqpoint{0.000000in}{-0.048611in}}%
\pgfusepath{stroke,fill}%
}%
\begin{pgfscope}%
\pgfsys@transformshift{2.616558in}{1.415315in}%
\pgfsys@useobject{currentmarker}{}%
\end{pgfscope}%
\end{pgfscope}%
\begin{pgfscope}%
\pgfpathrectangle{\pgfqpoint{0.514740in}{1.415315in}}{\pgfqpoint{2.720000in}{0.807339in}}%
\pgfusepath{clip}%
\pgfsetbuttcap%
\pgfsetroundjoin%
\pgfsetlinewidth{0.501875pt}%
\definecolor{currentstroke}{rgb}{0.690196,0.690196,0.690196}%
\pgfsetstrokecolor{currentstroke}%
\pgfsetstrokeopacity{0.250000}%
\pgfsetdash{{1.850000pt}{0.800000pt}}{0.000000pt}%
\pgfpathmoveto{\pgfqpoint{3.111103in}{1.415315in}}%
\pgfpathlineto{\pgfqpoint{3.111103in}{2.222654in}}%
\pgfusepath{stroke}%
\end{pgfscope}%
\begin{pgfscope}%
\pgfsetbuttcap%
\pgfsetroundjoin%
\definecolor{currentfill}{rgb}{0.000000,0.000000,0.000000}%
\pgfsetfillcolor{currentfill}%
\pgfsetlinewidth{0.803000pt}%
\definecolor{currentstroke}{rgb}{0.000000,0.000000,0.000000}%
\pgfsetstrokecolor{currentstroke}%
\pgfsetdash{}{0pt}%
\pgfsys@defobject{currentmarker}{\pgfqpoint{0.000000in}{-0.048611in}}{\pgfqpoint{0.000000in}{0.000000in}}{%
\pgfpathmoveto{\pgfqpoint{0.000000in}{0.000000in}}%
\pgfpathlineto{\pgfqpoint{0.000000in}{-0.048611in}}%
\pgfusepath{stroke,fill}%
}%
\begin{pgfscope}%
\pgfsys@transformshift{3.111103in}{1.415315in}%
\pgfsys@useobject{currentmarker}{}%
\end{pgfscope}%
\end{pgfscope}%
\begin{pgfscope}%
\pgfpathrectangle{\pgfqpoint{0.514740in}{1.415315in}}{\pgfqpoint{2.720000in}{0.807339in}}%
\pgfusepath{clip}%
\pgfsetbuttcap%
\pgfsetroundjoin%
\pgfsetlinewidth{0.501875pt}%
\definecolor{currentstroke}{rgb}{0.690196,0.690196,0.690196}%
\pgfsetstrokecolor{currentstroke}%
\pgfsetstrokeopacity{0.250000}%
\pgfsetdash{{1.850000pt}{0.800000pt}}{0.000000pt}%
\pgfpathmoveto{\pgfqpoint{0.514740in}{1.502768in}}%
\pgfpathlineto{\pgfqpoint{3.234740in}{1.502768in}}%
\pgfusepath{stroke}%
\end{pgfscope}%
\begin{pgfscope}%
\pgfsetbuttcap%
\pgfsetroundjoin%
\definecolor{currentfill}{rgb}{0.000000,0.000000,0.000000}%
\pgfsetfillcolor{currentfill}%
\pgfsetlinewidth{0.803000pt}%
\definecolor{currentstroke}{rgb}{0.000000,0.000000,0.000000}%
\pgfsetstrokecolor{currentstroke}%
\pgfsetdash{}{0pt}%
\pgfsys@defobject{currentmarker}{\pgfqpoint{-0.048611in}{0.000000in}}{\pgfqpoint{-0.000000in}{0.000000in}}{%
\pgfpathmoveto{\pgfqpoint{-0.000000in}{0.000000in}}%
\pgfpathlineto{\pgfqpoint{-0.048611in}{0.000000in}}%
\pgfusepath{stroke,fill}%
}%
\begin{pgfscope}%
\pgfsys@transformshift{0.514740in}{1.502768in}%
\pgfsys@useobject{currentmarker}{}%
\end{pgfscope}%
\end{pgfscope}%
\begin{pgfscope}%
\definecolor{textcolor}{rgb}{0.000000,0.000000,0.000000}%
\pgfsetstrokecolor{textcolor}%
\pgfsetfillcolor{textcolor}%
\pgftext[x=0.266667in, y=1.464188in, left, base]{\color{textcolor}{\rmfamily\fontsize{8.000000}{9.600000}\selectfont\catcode`\^=\active\def^{\ifmmode\sp\else\^{}\fi}\catcode`\%=\active\def%{\%}$\mathdefault{4.5}$}}%
\end{pgfscope}%
\begin{pgfscope}%
\pgfpathrectangle{\pgfqpoint{0.514740in}{1.415315in}}{\pgfqpoint{2.720000in}{0.807339in}}%
\pgfusepath{clip}%
\pgfsetbuttcap%
\pgfsetroundjoin%
\pgfsetlinewidth{0.501875pt}%
\definecolor{currentstroke}{rgb}{0.690196,0.690196,0.690196}%
\pgfsetstrokecolor{currentstroke}%
\pgfsetstrokeopacity{0.250000}%
\pgfsetdash{{1.850000pt}{0.800000pt}}{0.000000pt}%
\pgfpathmoveto{\pgfqpoint{0.514740in}{1.809367in}}%
\pgfpathlineto{\pgfqpoint{3.234740in}{1.809367in}}%
\pgfusepath{stroke}%
\end{pgfscope}%
\begin{pgfscope}%
\pgfsetbuttcap%
\pgfsetroundjoin%
\definecolor{currentfill}{rgb}{0.000000,0.000000,0.000000}%
\pgfsetfillcolor{currentfill}%
\pgfsetlinewidth{0.803000pt}%
\definecolor{currentstroke}{rgb}{0.000000,0.000000,0.000000}%
\pgfsetstrokecolor{currentstroke}%
\pgfsetdash{}{0pt}%
\pgfsys@defobject{currentmarker}{\pgfqpoint{-0.048611in}{0.000000in}}{\pgfqpoint{-0.000000in}{0.000000in}}{%
\pgfpathmoveto{\pgfqpoint{-0.000000in}{0.000000in}}%
\pgfpathlineto{\pgfqpoint{-0.048611in}{0.000000in}}%
\pgfusepath{stroke,fill}%
}%
\begin{pgfscope}%
\pgfsys@transformshift{0.514740in}{1.809367in}%
\pgfsys@useobject{currentmarker}{}%
\end{pgfscope}%
\end{pgfscope}%
\begin{pgfscope}%
\definecolor{textcolor}{rgb}{0.000000,0.000000,0.000000}%
\pgfsetstrokecolor{textcolor}%
\pgfsetfillcolor{textcolor}%
\pgftext[x=0.266667in, y=1.770787in, left, base]{\color{textcolor}{\rmfamily\fontsize{8.000000}{9.600000}\selectfont\catcode`\^=\active\def^{\ifmmode\sp\else\^{}\fi}\catcode`\%=\active\def%{\%}$\mathdefault{5.0}$}}%
\end{pgfscope}%
\begin{pgfscope}%
\pgfpathrectangle{\pgfqpoint{0.514740in}{1.415315in}}{\pgfqpoint{2.720000in}{0.807339in}}%
\pgfusepath{clip}%
\pgfsetbuttcap%
\pgfsetroundjoin%
\pgfsetlinewidth{0.501875pt}%
\definecolor{currentstroke}{rgb}{0.690196,0.690196,0.690196}%
\pgfsetstrokecolor{currentstroke}%
\pgfsetstrokeopacity{0.250000}%
\pgfsetdash{{1.850000pt}{0.800000pt}}{0.000000pt}%
\pgfpathmoveto{\pgfqpoint{0.514740in}{2.115966in}}%
\pgfpathlineto{\pgfqpoint{3.234740in}{2.115966in}}%
\pgfusepath{stroke}%
\end{pgfscope}%
\begin{pgfscope}%
\pgfsetbuttcap%
\pgfsetroundjoin%
\definecolor{currentfill}{rgb}{0.000000,0.000000,0.000000}%
\pgfsetfillcolor{currentfill}%
\pgfsetlinewidth{0.803000pt}%
\definecolor{currentstroke}{rgb}{0.000000,0.000000,0.000000}%
\pgfsetstrokecolor{currentstroke}%
\pgfsetdash{}{0pt}%
\pgfsys@defobject{currentmarker}{\pgfqpoint{-0.048611in}{0.000000in}}{\pgfqpoint{-0.000000in}{0.000000in}}{%
\pgfpathmoveto{\pgfqpoint{-0.000000in}{0.000000in}}%
\pgfpathlineto{\pgfqpoint{-0.048611in}{0.000000in}}%
\pgfusepath{stroke,fill}%
}%
\begin{pgfscope}%
\pgfsys@transformshift{0.514740in}{2.115966in}%
\pgfsys@useobject{currentmarker}{}%
\end{pgfscope}%
\end{pgfscope}%
\begin{pgfscope}%
\definecolor{textcolor}{rgb}{0.000000,0.000000,0.000000}%
\pgfsetstrokecolor{textcolor}%
\pgfsetfillcolor{textcolor}%
\pgftext[x=0.266667in, y=2.077386in, left, base]{\color{textcolor}{\rmfamily\fontsize{8.000000}{9.600000}\selectfont\catcode`\^=\active\def^{\ifmmode\sp\else\^{}\fi}\catcode`\%=\active\def%{\%}$\mathdefault{5.5}$}}%
\end{pgfscope}%
\begin{pgfscope}%
\definecolor{textcolor}{rgb}{0.000000,0.000000,0.000000}%
\pgfsetstrokecolor{textcolor}%
\pgfsetfillcolor{textcolor}%
\pgftext[x=0.211111in,y=1.818985in,,bottom,rotate=90.000000]{\color{textcolor}{\rmfamily\fontsize{8.000000}{9.600000}\selectfont\catcode`\^=\active\def^{\ifmmode\sp\else\^{}\fi}\catcode`\%=\active\def%{\%}$v_o$ (V)}}%
\end{pgfscope}%
\begin{pgfscope}%
\pgfpathrectangle{\pgfqpoint{0.514740in}{1.415315in}}{\pgfqpoint{2.720000in}{0.807339in}}%
\pgfusepath{clip}%
\pgfsetrectcap%
\pgfsetroundjoin%
\pgfsetlinewidth{1.505625pt}%
\definecolor{currentstroke}{rgb}{0.000000,0.447059,0.698039}%
\pgfsetstrokecolor{currentstroke}%
\pgfsetdash{}{0pt}%
\pgfpathmoveto{\pgfqpoint{0.638376in}{1.809362in}}%
\pgfpathlineto{\pgfqpoint{0.647666in}{1.766187in}}%
\pgfpathlineto{\pgfqpoint{0.657246in}{1.727767in}}%
\pgfpathlineto{\pgfqpoint{0.667137in}{1.694546in}}%
\pgfpathlineto{\pgfqpoint{0.677028in}{1.667801in}}%
\pgfpathlineto{\pgfqpoint{0.686919in}{1.647454in}}%
\pgfpathlineto{\pgfqpoint{0.691865in}{1.639655in}}%
\pgfpathlineto{\pgfqpoint{0.696810in}{1.633424in}}%
\pgfpathlineto{\pgfqpoint{0.701756in}{1.628751in}}%
\pgfpathlineto{\pgfqpoint{0.706701in}{1.625623in}}%
\pgfpathlineto{\pgfqpoint{0.711647in}{1.624031in}}%
\pgfpathlineto{\pgfqpoint{0.716592in}{1.623960in}}%
\pgfpathlineto{\pgfqpoint{0.721537in}{1.625401in}}%
\pgfpathlineto{\pgfqpoint{0.726483in}{1.628340in}}%
\pgfpathlineto{\pgfqpoint{0.731428in}{1.632764in}}%
\pgfpathlineto{\pgfqpoint{0.736374in}{1.638661in}}%
\pgfpathlineto{\pgfqpoint{0.741319in}{1.646018in}}%
\pgfpathlineto{\pgfqpoint{0.751210in}{1.665056in}}%
\pgfpathlineto{\pgfqpoint{0.761101in}{1.689768in}}%
\pgfpathlineto{\pgfqpoint{0.770992in}{1.720042in}}%
\pgfpathlineto{\pgfqpoint{0.780883in}{1.755759in}}%
\pgfpathlineto{\pgfqpoint{0.790774in}{1.796799in}}%
\pgfpathlineto{\pgfqpoint{0.804518in}{1.862354in}}%
\pgfpathlineto{\pgfqpoint{0.818207in}{1.927223in}}%
\pgfpathlineto{\pgfqpoint{0.833043in}{1.989120in}}%
\pgfpathlineto{\pgfqpoint{0.847880in}{2.042345in}}%
\pgfpathlineto{\pgfqpoint{0.857771in}{2.073057in}}%
\pgfpathlineto{\pgfqpoint{0.867661in}{2.099988in}}%
\pgfpathlineto{\pgfqpoint{0.877552in}{2.123173in}}%
\pgfpathlineto{\pgfqpoint{0.887443in}{2.142651in}}%
\pgfpathlineto{\pgfqpoint{0.897334in}{2.158462in}}%
\pgfpathlineto{\pgfqpoint{0.907225in}{2.170650in}}%
\pgfpathlineto{\pgfqpoint{0.917116in}{2.179262in}}%
\pgfpathlineto{\pgfqpoint{0.927007in}{2.184347in}}%
\pgfpathlineto{\pgfqpoint{0.931952in}{2.185583in}}%
\pgfpathlineto{\pgfqpoint{0.936898in}{2.185957in}}%
\pgfpathlineto{\pgfqpoint{0.941843in}{2.185476in}}%
\pgfpathlineto{\pgfqpoint{0.946789in}{2.184147in}}%
\pgfpathlineto{\pgfqpoint{0.956680in}{2.178974in}}%
\pgfpathlineto{\pgfqpoint{0.966571in}{2.170497in}}%
\pgfpathlineto{\pgfqpoint{0.976461in}{2.158777in}}%
\pgfpathlineto{\pgfqpoint{0.986352in}{2.143879in}}%
\pgfpathlineto{\pgfqpoint{0.996243in}{2.125868in}}%
\pgfpathlineto{\pgfqpoint{1.006134in}{2.104813in}}%
\pgfpathlineto{\pgfqpoint{1.016025in}{2.080783in}}%
\pgfpathlineto{\pgfqpoint{1.030862in}{2.039318in}}%
\pgfpathlineto{\pgfqpoint{1.045698in}{1.991572in}}%
\pgfpathlineto{\pgfqpoint{1.060534in}{1.937804in}}%
\pgfpathlineto{\pgfqpoint{1.075371in}{1.878281in}}%
\pgfpathlineto{\pgfqpoint{1.090207in}{1.813283in}}%
\pgfpathlineto{\pgfqpoint{1.112060in}{1.709340in}}%
\pgfpathlineto{\pgfqpoint{1.138289in}{1.586942in}}%
\pgfpathlineto{\pgfqpoint{1.148180in}{1.548076in}}%
\pgfpathlineto{\pgfqpoint{1.158071in}{1.515915in}}%
\pgfpathlineto{\pgfqpoint{1.167962in}{1.490383in}}%
\pgfpathlineto{\pgfqpoint{1.177853in}{1.471399in}}%
\pgfpathlineto{\pgfqpoint{1.182798in}{1.464336in}}%
\pgfpathlineto{\pgfqpoint{1.187744in}{1.458877in}}%
\pgfpathlineto{\pgfqpoint{1.192689in}{1.455012in}}%
\pgfpathlineto{\pgfqpoint{1.197635in}{1.452727in}}%
\pgfpathlineto{\pgfqpoint{1.202580in}{1.452012in}}%
\pgfpathlineto{\pgfqpoint{1.207526in}{1.452854in}}%
\pgfpathlineto{\pgfqpoint{1.212471in}{1.455240in}}%
\pgfpathlineto{\pgfqpoint{1.217417in}{1.459157in}}%
\pgfpathlineto{\pgfqpoint{1.222362in}{1.464593in}}%
\pgfpathlineto{\pgfqpoint{1.227308in}{1.471533in}}%
\pgfpathlineto{\pgfqpoint{1.237198in}{1.489873in}}%
\pgfpathlineto{\pgfqpoint{1.247089in}{1.514066in}}%
\pgfpathlineto{\pgfqpoint{1.256980in}{1.543995in}}%
\pgfpathlineto{\pgfqpoint{1.266871in}{1.579541in}}%
\pgfpathlineto{\pgfqpoint{1.276762in}{1.620581in}}%
\pgfpathlineto{\pgfqpoint{1.286653in}{1.666988in}}%
\pgfpathlineto{\pgfqpoint{1.302862in}{1.753109in}}%
\pgfpathlineto{\pgfqpoint{1.317698in}{1.828035in}}%
\pgfpathlineto{\pgfqpoint{1.332534in}{1.894322in}}%
\pgfpathlineto{\pgfqpoint{1.347371in}{1.952036in}}%
\pgfpathlineto{\pgfqpoint{1.357262in}{1.985788in}}%
\pgfpathlineto{\pgfqpoint{1.367153in}{2.015794in}}%
\pgfpathlineto{\pgfqpoint{1.377043in}{2.042085in}}%
\pgfpathlineto{\pgfqpoint{1.386934in}{2.064696in}}%
\pgfpathlineto{\pgfqpoint{1.396825in}{2.083665in}}%
\pgfpathlineto{\pgfqpoint{1.406716in}{2.099031in}}%
\pgfpathlineto{\pgfqpoint{1.416607in}{2.110839in}}%
\pgfpathlineto{\pgfqpoint{1.426498in}{2.119134in}}%
\pgfpathlineto{\pgfqpoint{1.436389in}{2.123964in}}%
\pgfpathlineto{\pgfqpoint{1.441334in}{2.125096in}}%
\pgfpathlineto{\pgfqpoint{1.446280in}{2.125381in}}%
\pgfpathlineto{\pgfqpoint{1.451225in}{2.124827in}}%
\pgfpathlineto{\pgfqpoint{1.456171in}{2.123440in}}%
\pgfpathlineto{\pgfqpoint{1.466062in}{2.118195in}}%
\pgfpathlineto{\pgfqpoint{1.475953in}{2.109706in}}%
\pgfpathlineto{\pgfqpoint{1.485844in}{2.098032in}}%
\pgfpathlineto{\pgfqpoint{1.495734in}{2.083238in}}%
\pgfpathlineto{\pgfqpoint{1.505625in}{2.065387in}}%
\pgfpathlineto{\pgfqpoint{1.515516in}{2.044547in}}%
\pgfpathlineto{\pgfqpoint{1.525407in}{2.020787in}}%
\pgfpathlineto{\pgfqpoint{1.540244in}{1.979826in}}%
\pgfpathlineto{\pgfqpoint{1.555080in}{1.932700in}}%
\pgfpathlineto{\pgfqpoint{1.569916in}{1.879664in}}%
\pgfpathlineto{\pgfqpoint{1.584753in}{1.820981in}}%
\pgfpathlineto{\pgfqpoint{1.599589in}{1.756926in}}%
\pgfpathlineto{\pgfqpoint{1.624865in}{1.638605in}}%
\pgfpathlineto{\pgfqpoint{1.633800in}{1.598256in}}%
\pgfpathlineto{\pgfqpoint{1.643691in}{1.559781in}}%
\pgfpathlineto{\pgfqpoint{1.653582in}{1.527994in}}%
\pgfpathlineto{\pgfqpoint{1.663473in}{1.502818in}}%
\pgfpathlineto{\pgfqpoint{1.673364in}{1.484173in}}%
\pgfpathlineto{\pgfqpoint{1.678309in}{1.477273in}}%
\pgfpathlineto{\pgfqpoint{1.683255in}{1.471973in}}%
\pgfpathlineto{\pgfqpoint{1.688200in}{1.468261in}}%
\pgfpathlineto{\pgfqpoint{1.693146in}{1.466127in}}%
\pgfpathlineto{\pgfqpoint{1.698091in}{1.465557in}}%
\pgfpathlineto{\pgfqpoint{1.703037in}{1.466539in}}%
\pgfpathlineto{\pgfqpoint{1.707982in}{1.469062in}}%
\pgfpathlineto{\pgfqpoint{1.712928in}{1.473111in}}%
\pgfpathlineto{\pgfqpoint{1.717873in}{1.478674in}}%
\pgfpathlineto{\pgfqpoint{1.722818in}{1.485738in}}%
\pgfpathlineto{\pgfqpoint{1.732709in}{1.504312in}}%
\pgfpathlineto{\pgfqpoint{1.742600in}{1.528721in}}%
\pgfpathlineto{\pgfqpoint{1.752491in}{1.558848in}}%
\pgfpathlineto{\pgfqpoint{1.762382in}{1.594576in}}%
\pgfpathlineto{\pgfqpoint{1.772273in}{1.635780in}}%
\pgfpathlineto{\pgfqpoint{1.782164in}{1.682334in}}%
\pgfpathlineto{\pgfqpoint{1.817189in}{1.860282in}}%
\pgfpathlineto{\pgfqpoint{1.832025in}{1.923125in}}%
\pgfpathlineto{\pgfqpoint{1.846862in}{1.977415in}}%
\pgfpathlineto{\pgfqpoint{1.856753in}{2.008899in}}%
\pgfpathlineto{\pgfqpoint{1.866644in}{2.036652in}}%
\pgfpathlineto{\pgfqpoint{1.876535in}{2.060706in}}%
\pgfpathlineto{\pgfqpoint{1.886425in}{2.081099in}}%
\pgfpathlineto{\pgfqpoint{1.896316in}{2.097868in}}%
\pgfpathlineto{\pgfqpoint{1.906207in}{2.111058in}}%
\pgfpathlineto{\pgfqpoint{1.916098in}{2.120712in}}%
\pgfpathlineto{\pgfqpoint{1.925989in}{2.126878in}}%
\pgfpathlineto{\pgfqpoint{1.930935in}{2.128669in}}%
\pgfpathlineto{\pgfqpoint{1.935880in}{2.129608in}}%
\pgfpathlineto{\pgfqpoint{1.940826in}{2.129700in}}%
\pgfpathlineto{\pgfqpoint{1.945771in}{2.128952in}}%
\pgfpathlineto{\pgfqpoint{1.950716in}{2.127373in}}%
\pgfpathlineto{\pgfqpoint{1.960607in}{2.121745in}}%
\pgfpathlineto{\pgfqpoint{1.970498in}{2.112876in}}%
\pgfpathlineto{\pgfqpoint{1.980389in}{2.100826in}}%
\pgfpathlineto{\pgfqpoint{1.990280in}{2.085658in}}%
\pgfpathlineto{\pgfqpoint{2.000171in}{2.067438in}}%
\pgfpathlineto{\pgfqpoint{2.010062in}{2.046232in}}%
\pgfpathlineto{\pgfqpoint{2.019953in}{2.022111in}}%
\pgfpathlineto{\pgfqpoint{2.034789in}{1.980617in}}%
\pgfpathlineto{\pgfqpoint{2.049626in}{1.932969in}}%
\pgfpathlineto{\pgfqpoint{2.064462in}{1.879424in}}%
\pgfpathlineto{\pgfqpoint{2.079298in}{1.820246in}}%
\pgfpathlineto{\pgfqpoint{2.094135in}{1.755710in}}%
\pgfpathlineto{\pgfqpoint{2.133289in}{1.576652in}}%
\pgfpathlineto{\pgfqpoint{2.143180in}{1.541555in}}%
\pgfpathlineto{\pgfqpoint{2.153071in}{1.513110in}}%
\pgfpathlineto{\pgfqpoint{2.162962in}{1.491238in}}%
\pgfpathlineto{\pgfqpoint{2.167908in}{1.482740in}}%
\pgfpathlineto{\pgfqpoint{2.172853in}{1.475854in}}%
\pgfpathlineto{\pgfqpoint{2.177799in}{1.470569in}}%
\pgfpathlineto{\pgfqpoint{2.182744in}{1.466872in}}%
\pgfpathlineto{\pgfqpoint{2.187689in}{1.464752in}}%
\pgfpathlineto{\pgfqpoint{2.192635in}{1.464198in}}%
\pgfpathlineto{\pgfqpoint{2.197580in}{1.465196in}}%
\pgfpathlineto{\pgfqpoint{2.202526in}{1.467734in}}%
\pgfpathlineto{\pgfqpoint{2.207471in}{1.471799in}}%
\pgfpathlineto{\pgfqpoint{2.212417in}{1.477378in}}%
\pgfpathlineto{\pgfqpoint{2.217362in}{1.484458in}}%
\pgfpathlineto{\pgfqpoint{2.227253in}{1.503064in}}%
\pgfpathlineto{\pgfqpoint{2.237144in}{1.527507in}}%
\pgfpathlineto{\pgfqpoint{2.247035in}{1.557669in}}%
\pgfpathlineto{\pgfqpoint{2.256926in}{1.593431in}}%
\pgfpathlineto{\pgfqpoint{2.266817in}{1.634670in}}%
\pgfpathlineto{\pgfqpoint{2.276708in}{1.681260in}}%
\pgfpathlineto{\pgfqpoint{2.314521in}{1.871841in}}%
\pgfpathlineto{\pgfqpoint{2.329358in}{1.933140in}}%
\pgfpathlineto{\pgfqpoint{2.344194in}{1.985901in}}%
\pgfpathlineto{\pgfqpoint{2.354085in}{2.016377in}}%
\pgfpathlineto{\pgfqpoint{2.363976in}{2.043130in}}%
\pgfpathlineto{\pgfqpoint{2.373867in}{2.066195in}}%
\pgfpathlineto{\pgfqpoint{2.383758in}{2.085608in}}%
\pgfpathlineto{\pgfqpoint{2.393648in}{2.101410in}}%
\pgfpathlineto{\pgfqpoint{2.403539in}{2.113644in}}%
\pgfpathlineto{\pgfqpoint{2.413430in}{2.122356in}}%
\pgfpathlineto{\pgfqpoint{2.423321in}{2.127593in}}%
\pgfpathlineto{\pgfqpoint{2.428267in}{2.128925in}}%
\pgfpathlineto{\pgfqpoint{2.433212in}{2.129408in}}%
\pgfpathlineto{\pgfqpoint{2.438158in}{2.129048in}}%
\pgfpathlineto{\pgfqpoint{2.443103in}{2.127853in}}%
\pgfpathlineto{\pgfqpoint{2.452994in}{2.122984in}}%
\pgfpathlineto{\pgfqpoint{2.462885in}{2.114860in}}%
\pgfpathlineto{\pgfqpoint{2.472776in}{2.103542in}}%
\pgfpathlineto{\pgfqpoint{2.482667in}{2.089091in}}%
\pgfpathlineto{\pgfqpoint{2.492558in}{2.071573in}}%
\pgfpathlineto{\pgfqpoint{2.502448in}{2.051054in}}%
\pgfpathlineto{\pgfqpoint{2.512339in}{2.027604in}}%
\pgfpathlineto{\pgfqpoint{2.527176in}{1.987085in}}%
\pgfpathlineto{\pgfqpoint{2.542012in}{1.940375in}}%
\pgfpathlineto{\pgfqpoint{2.556849in}{1.887728in}}%
\pgfpathlineto{\pgfqpoint{2.571685in}{1.829406in}}%
\pgfpathlineto{\pgfqpoint{2.586521in}{1.765685in}}%
\pgfpathlineto{\pgfqpoint{2.609175in}{1.659782in}}%
\pgfpathlineto{\pgfqpoint{2.621932in}{1.601064in}}%
\pgfpathlineto{\pgfqpoint{2.631823in}{1.561963in}}%
\pgfpathlineto{\pgfqpoint{2.641714in}{1.529558in}}%
\pgfpathlineto{\pgfqpoint{2.651605in}{1.503773in}}%
\pgfpathlineto{\pgfqpoint{2.661496in}{1.484527in}}%
\pgfpathlineto{\pgfqpoint{2.666441in}{1.477330in}}%
\pgfpathlineto{\pgfqpoint{2.671387in}{1.471735in}}%
\pgfpathlineto{\pgfqpoint{2.676332in}{1.467731in}}%
\pgfpathlineto{\pgfqpoint{2.681278in}{1.465307in}}%
\pgfpathlineto{\pgfqpoint{2.686223in}{1.464450in}}%
\pgfpathlineto{\pgfqpoint{2.691169in}{1.465148in}}%
\pgfpathlineto{\pgfqpoint{2.696114in}{1.467388in}}%
\pgfpathlineto{\pgfqpoint{2.701059in}{1.471158in}}%
\pgfpathlineto{\pgfqpoint{2.706005in}{1.476444in}}%
\pgfpathlineto{\pgfqpoint{2.710950in}{1.483234in}}%
\pgfpathlineto{\pgfqpoint{2.720841in}{1.501268in}}%
\pgfpathlineto{\pgfqpoint{2.730732in}{1.525149in}}%
\pgfpathlineto{\pgfqpoint{2.740623in}{1.554760in}}%
\pgfpathlineto{\pgfqpoint{2.750514in}{1.589983in}}%
\pgfpathlineto{\pgfqpoint{2.760405in}{1.630695in}}%
\pgfpathlineto{\pgfqpoint{2.770296in}{1.676770in}}%
\pgfpathlineto{\pgfqpoint{2.785360in}{1.756240in}}%
\pgfpathlineto{\pgfqpoint{2.799176in}{1.826436in}}%
\pgfpathlineto{\pgfqpoint{2.814012in}{1.893464in}}%
\pgfpathlineto{\pgfqpoint{2.828849in}{1.951903in}}%
\pgfpathlineto{\pgfqpoint{2.843685in}{2.001832in}}%
\pgfpathlineto{\pgfqpoint{2.853576in}{2.030439in}}%
\pgfpathlineto{\pgfqpoint{2.863467in}{2.055339in}}%
\pgfpathlineto{\pgfqpoint{2.873358in}{2.076568in}}%
\pgfpathlineto{\pgfqpoint{2.883249in}{2.094166in}}%
\pgfpathlineto{\pgfqpoint{2.893140in}{2.108174in}}%
\pgfpathlineto{\pgfqpoint{2.903030in}{2.118636in}}%
\pgfpathlineto{\pgfqpoint{2.912921in}{2.125600in}}%
\pgfpathlineto{\pgfqpoint{2.917867in}{2.127785in}}%
\pgfpathlineto{\pgfqpoint{2.922812in}{2.129114in}}%
\pgfpathlineto{\pgfqpoint{2.927758in}{2.129594in}}%
\pgfpathlineto{\pgfqpoint{2.932703in}{2.129232in}}%
\pgfpathlineto{\pgfqpoint{2.937649in}{2.128034in}}%
\pgfpathlineto{\pgfqpoint{2.947540in}{2.123160in}}%
\pgfpathlineto{\pgfqpoint{2.957430in}{2.115031in}}%
\pgfpathlineto{\pgfqpoint{2.967321in}{2.103707in}}%
\pgfpathlineto{\pgfqpoint{2.977212in}{2.089251in}}%
\pgfpathlineto{\pgfqpoint{2.987103in}{2.071727in}}%
\pgfpathlineto{\pgfqpoint{2.996994in}{2.051203in}}%
\pgfpathlineto{\pgfqpoint{3.006885in}{2.027747in}}%
\pgfpathlineto{\pgfqpoint{3.021721in}{1.987220in}}%
\pgfpathlineto{\pgfqpoint{3.036558in}{1.940501in}}%
\pgfpathlineto{\pgfqpoint{3.051394in}{1.887845in}}%
\pgfpathlineto{\pgfqpoint{3.066231in}{1.829515in}}%
\pgfpathlineto{\pgfqpoint{3.081067in}{1.765785in}}%
\pgfpathlineto{\pgfqpoint{3.103705in}{1.659947in}}%
\pgfpathlineto{\pgfqpoint{3.111103in}{1.625231in}}%
\pgfpathlineto{\pgfqpoint{3.111103in}{1.625231in}}%
\pgfusepath{stroke}%
\end{pgfscope}%
\begin{pgfscope}%
\pgfsetrectcap%
\pgfsetmiterjoin%
\pgfsetlinewidth{0.803000pt}%
\definecolor{currentstroke}{rgb}{0.000000,0.000000,0.000000}%
\pgfsetstrokecolor{currentstroke}%
\pgfsetdash{}{0pt}%
\pgfpathmoveto{\pgfqpoint{0.514740in}{1.415315in}}%
\pgfpathlineto{\pgfqpoint{0.514740in}{2.222654in}}%
\pgfusepath{stroke}%
\end{pgfscope}%
\begin{pgfscope}%
\pgfsetrectcap%
\pgfsetmiterjoin%
\pgfsetlinewidth{0.803000pt}%
\definecolor{currentstroke}{rgb}{0.000000,0.000000,0.000000}%
\pgfsetstrokecolor{currentstroke}%
\pgfsetdash{}{0pt}%
\pgfpathmoveto{\pgfqpoint{3.234740in}{1.415315in}}%
\pgfpathlineto{\pgfqpoint{3.234740in}{2.222654in}}%
\pgfusepath{stroke}%
\end{pgfscope}%
\begin{pgfscope}%
\pgfsetrectcap%
\pgfsetmiterjoin%
\pgfsetlinewidth{0.803000pt}%
\definecolor{currentstroke}{rgb}{0.000000,0.000000,0.000000}%
\pgfsetstrokecolor{currentstroke}%
\pgfsetdash{}{0pt}%
\pgfpathmoveto{\pgfqpoint{0.514740in}{1.415315in}}%
\pgfpathlineto{\pgfqpoint{3.234740in}{1.415315in}}%
\pgfusepath{stroke}%
\end{pgfscope}%
\begin{pgfscope}%
\pgfsetrectcap%
\pgfsetmiterjoin%
\pgfsetlinewidth{0.803000pt}%
\definecolor{currentstroke}{rgb}{0.000000,0.000000,0.000000}%
\pgfsetstrokecolor{currentstroke}%
\pgfsetdash{}{0pt}%
\pgfpathmoveto{\pgfqpoint{0.514740in}{2.222654in}}%
\pgfpathlineto{\pgfqpoint{3.234740in}{2.222654in}}%
\pgfusepath{stroke}%
\end{pgfscope}%
\begin{pgfscope}%
\definecolor{textcolor}{rgb}{0.000000,0.000000,0.000000}%
\pgfsetstrokecolor{textcolor}%
\pgfsetfillcolor{textcolor}%
\pgftext[x=0.569140in,y=2.166141in,left,top]{\color{textcolor}{\rmfamily\fontsize{8.000000}{9.600000}\selectfont\catcode`\^=\active\def^{\ifmmode\sp\else\^{}\fi}\catcode`\%=\active\def%{\%}(a)}}%
\end{pgfscope}%
\begin{pgfscope}%
\pgfsetbuttcap%
\pgfsetmiterjoin%
\definecolor{currentfill}{rgb}{1.000000,1.000000,1.000000}%
\pgfsetfillcolor{currentfill}%
\pgfsetlinewidth{0.000000pt}%
\definecolor{currentstroke}{rgb}{0.000000,0.000000,0.000000}%
\pgfsetstrokecolor{currentstroke}%
\pgfsetstrokeopacity{0.000000}%
\pgfsetdash{}{0pt}%
\pgfpathmoveto{\pgfqpoint{0.514740in}{0.462654in}}%
\pgfpathlineto{\pgfqpoint{3.234740in}{0.462654in}}%
\pgfpathlineto{\pgfqpoint{3.234740in}{1.269994in}}%
\pgfpathlineto{\pgfqpoint{0.514740in}{1.269994in}}%
\pgfpathlineto{\pgfqpoint{0.514740in}{0.462654in}}%
\pgfpathclose%
\pgfusepath{fill}%
\end{pgfscope}%
\begin{pgfscope}%
\pgfpathrectangle{\pgfqpoint{0.514740in}{0.462654in}}{\pgfqpoint{2.720000in}{0.807339in}}%
\pgfusepath{clip}%
\pgfsetbuttcap%
\pgfsetroundjoin%
\pgfsetlinewidth{0.501875pt}%
\definecolor{currentstroke}{rgb}{0.690196,0.690196,0.690196}%
\pgfsetstrokecolor{currentstroke}%
\pgfsetstrokeopacity{0.250000}%
\pgfsetdash{{1.850000pt}{0.800000pt}}{0.000000pt}%
\pgfpathmoveto{\pgfqpoint{0.638375in}{0.462654in}}%
\pgfpathlineto{\pgfqpoint{0.638375in}{1.269994in}}%
\pgfusepath{stroke}%
\end{pgfscope}%
\begin{pgfscope}%
\pgfsetbuttcap%
\pgfsetroundjoin%
\definecolor{currentfill}{rgb}{0.000000,0.000000,0.000000}%
\pgfsetfillcolor{currentfill}%
\pgfsetlinewidth{0.803000pt}%
\definecolor{currentstroke}{rgb}{0.000000,0.000000,0.000000}%
\pgfsetstrokecolor{currentstroke}%
\pgfsetdash{}{0pt}%
\pgfsys@defobject{currentmarker}{\pgfqpoint{0.000000in}{-0.048611in}}{\pgfqpoint{0.000000in}{0.000000in}}{%
\pgfpathmoveto{\pgfqpoint{0.000000in}{0.000000in}}%
\pgfpathlineto{\pgfqpoint{0.000000in}{-0.048611in}}%
\pgfusepath{stroke,fill}%
}%
\begin{pgfscope}%
\pgfsys@transformshift{0.638375in}{0.462654in}%
\pgfsys@useobject{currentmarker}{}%
\end{pgfscope}%
\end{pgfscope}%
\begin{pgfscope}%
\definecolor{textcolor}{rgb}{0.000000,0.000000,0.000000}%
\pgfsetstrokecolor{textcolor}%
\pgfsetfillcolor{textcolor}%
\pgftext[x=0.638375in,y=0.365432in,,top]{\color{textcolor}{\rmfamily\fontsize{8.000000}{9.600000}\selectfont\catcode`\^=\active\def^{\ifmmode\sp\else\^{}\fi}\catcode`\%=\active\def%{\%}$\mathdefault{0}$}}%
\end{pgfscope}%
\begin{pgfscope}%
\pgfpathrectangle{\pgfqpoint{0.514740in}{0.462654in}}{\pgfqpoint{2.720000in}{0.807339in}}%
\pgfusepath{clip}%
\pgfsetbuttcap%
\pgfsetroundjoin%
\pgfsetlinewidth{0.501875pt}%
\definecolor{currentstroke}{rgb}{0.690196,0.690196,0.690196}%
\pgfsetstrokecolor{currentstroke}%
\pgfsetstrokeopacity{0.250000}%
\pgfsetdash{{1.850000pt}{0.800000pt}}{0.000000pt}%
\pgfpathmoveto{\pgfqpoint{1.132921in}{0.462654in}}%
\pgfpathlineto{\pgfqpoint{1.132921in}{1.269994in}}%
\pgfusepath{stroke}%
\end{pgfscope}%
\begin{pgfscope}%
\pgfsetbuttcap%
\pgfsetroundjoin%
\definecolor{currentfill}{rgb}{0.000000,0.000000,0.000000}%
\pgfsetfillcolor{currentfill}%
\pgfsetlinewidth{0.803000pt}%
\definecolor{currentstroke}{rgb}{0.000000,0.000000,0.000000}%
\pgfsetstrokecolor{currentstroke}%
\pgfsetdash{}{0pt}%
\pgfsys@defobject{currentmarker}{\pgfqpoint{0.000000in}{-0.048611in}}{\pgfqpoint{0.000000in}{0.000000in}}{%
\pgfpathmoveto{\pgfqpoint{0.000000in}{0.000000in}}%
\pgfpathlineto{\pgfqpoint{0.000000in}{-0.048611in}}%
\pgfusepath{stroke,fill}%
}%
\begin{pgfscope}%
\pgfsys@transformshift{1.132921in}{0.462654in}%
\pgfsys@useobject{currentmarker}{}%
\end{pgfscope}%
\end{pgfscope}%
\begin{pgfscope}%
\definecolor{textcolor}{rgb}{0.000000,0.000000,0.000000}%
\pgfsetstrokecolor{textcolor}%
\pgfsetfillcolor{textcolor}%
\pgftext[x=1.132921in,y=0.365432in,,top]{\color{textcolor}{\rmfamily\fontsize{8.000000}{9.600000}\selectfont\catcode`\^=\active\def^{\ifmmode\sp\else\^{}\fi}\catcode`\%=\active\def%{\%}$\mathdefault{5}$}}%
\end{pgfscope}%
\begin{pgfscope}%
\pgfpathrectangle{\pgfqpoint{0.514740in}{0.462654in}}{\pgfqpoint{2.720000in}{0.807339in}}%
\pgfusepath{clip}%
\pgfsetbuttcap%
\pgfsetroundjoin%
\pgfsetlinewidth{0.501875pt}%
\definecolor{currentstroke}{rgb}{0.690196,0.690196,0.690196}%
\pgfsetstrokecolor{currentstroke}%
\pgfsetstrokeopacity{0.250000}%
\pgfsetdash{{1.850000pt}{0.800000pt}}{0.000000pt}%
\pgfpathmoveto{\pgfqpoint{1.627467in}{0.462654in}}%
\pgfpathlineto{\pgfqpoint{1.627467in}{1.269994in}}%
\pgfusepath{stroke}%
\end{pgfscope}%
\begin{pgfscope}%
\pgfsetbuttcap%
\pgfsetroundjoin%
\definecolor{currentfill}{rgb}{0.000000,0.000000,0.000000}%
\pgfsetfillcolor{currentfill}%
\pgfsetlinewidth{0.803000pt}%
\definecolor{currentstroke}{rgb}{0.000000,0.000000,0.000000}%
\pgfsetstrokecolor{currentstroke}%
\pgfsetdash{}{0pt}%
\pgfsys@defobject{currentmarker}{\pgfqpoint{0.000000in}{-0.048611in}}{\pgfqpoint{0.000000in}{0.000000in}}{%
\pgfpathmoveto{\pgfqpoint{0.000000in}{0.000000in}}%
\pgfpathlineto{\pgfqpoint{0.000000in}{-0.048611in}}%
\pgfusepath{stroke,fill}%
}%
\begin{pgfscope}%
\pgfsys@transformshift{1.627467in}{0.462654in}%
\pgfsys@useobject{currentmarker}{}%
\end{pgfscope}%
\end{pgfscope}%
\begin{pgfscope}%
\definecolor{textcolor}{rgb}{0.000000,0.000000,0.000000}%
\pgfsetstrokecolor{textcolor}%
\pgfsetfillcolor{textcolor}%
\pgftext[x=1.627467in,y=0.365432in,,top]{\color{textcolor}{\rmfamily\fontsize{8.000000}{9.600000}\selectfont\catcode`\^=\active\def^{\ifmmode\sp\else\^{}\fi}\catcode`\%=\active\def%{\%}$\mathdefault{10}$}}%
\end{pgfscope}%
\begin{pgfscope}%
\pgfpathrectangle{\pgfqpoint{0.514740in}{0.462654in}}{\pgfqpoint{2.720000in}{0.807339in}}%
\pgfusepath{clip}%
\pgfsetbuttcap%
\pgfsetroundjoin%
\pgfsetlinewidth{0.501875pt}%
\definecolor{currentstroke}{rgb}{0.690196,0.690196,0.690196}%
\pgfsetstrokecolor{currentstroke}%
\pgfsetstrokeopacity{0.250000}%
\pgfsetdash{{1.850000pt}{0.800000pt}}{0.000000pt}%
\pgfpathmoveto{\pgfqpoint{2.122012in}{0.462654in}}%
\pgfpathlineto{\pgfqpoint{2.122012in}{1.269994in}}%
\pgfusepath{stroke}%
\end{pgfscope}%
\begin{pgfscope}%
\pgfsetbuttcap%
\pgfsetroundjoin%
\definecolor{currentfill}{rgb}{0.000000,0.000000,0.000000}%
\pgfsetfillcolor{currentfill}%
\pgfsetlinewidth{0.803000pt}%
\definecolor{currentstroke}{rgb}{0.000000,0.000000,0.000000}%
\pgfsetstrokecolor{currentstroke}%
\pgfsetdash{}{0pt}%
\pgfsys@defobject{currentmarker}{\pgfqpoint{0.000000in}{-0.048611in}}{\pgfqpoint{0.000000in}{0.000000in}}{%
\pgfpathmoveto{\pgfqpoint{0.000000in}{0.000000in}}%
\pgfpathlineto{\pgfqpoint{0.000000in}{-0.048611in}}%
\pgfusepath{stroke,fill}%
}%
\begin{pgfscope}%
\pgfsys@transformshift{2.122012in}{0.462654in}%
\pgfsys@useobject{currentmarker}{}%
\end{pgfscope}%
\end{pgfscope}%
\begin{pgfscope}%
\definecolor{textcolor}{rgb}{0.000000,0.000000,0.000000}%
\pgfsetstrokecolor{textcolor}%
\pgfsetfillcolor{textcolor}%
\pgftext[x=2.122012in,y=0.365432in,,top]{\color{textcolor}{\rmfamily\fontsize{8.000000}{9.600000}\selectfont\catcode`\^=\active\def^{\ifmmode\sp\else\^{}\fi}\catcode`\%=\active\def%{\%}$\mathdefault{15}$}}%
\end{pgfscope}%
\begin{pgfscope}%
\pgfpathrectangle{\pgfqpoint{0.514740in}{0.462654in}}{\pgfqpoint{2.720000in}{0.807339in}}%
\pgfusepath{clip}%
\pgfsetbuttcap%
\pgfsetroundjoin%
\pgfsetlinewidth{0.501875pt}%
\definecolor{currentstroke}{rgb}{0.690196,0.690196,0.690196}%
\pgfsetstrokecolor{currentstroke}%
\pgfsetstrokeopacity{0.250000}%
\pgfsetdash{{1.850000pt}{0.800000pt}}{0.000000pt}%
\pgfpathmoveto{\pgfqpoint{2.616558in}{0.462654in}}%
\pgfpathlineto{\pgfqpoint{2.616558in}{1.269994in}}%
\pgfusepath{stroke}%
\end{pgfscope}%
\begin{pgfscope}%
\pgfsetbuttcap%
\pgfsetroundjoin%
\definecolor{currentfill}{rgb}{0.000000,0.000000,0.000000}%
\pgfsetfillcolor{currentfill}%
\pgfsetlinewidth{0.803000pt}%
\definecolor{currentstroke}{rgb}{0.000000,0.000000,0.000000}%
\pgfsetstrokecolor{currentstroke}%
\pgfsetdash{}{0pt}%
\pgfsys@defobject{currentmarker}{\pgfqpoint{0.000000in}{-0.048611in}}{\pgfqpoint{0.000000in}{0.000000in}}{%
\pgfpathmoveto{\pgfqpoint{0.000000in}{0.000000in}}%
\pgfpathlineto{\pgfqpoint{0.000000in}{-0.048611in}}%
\pgfusepath{stroke,fill}%
}%
\begin{pgfscope}%
\pgfsys@transformshift{2.616558in}{0.462654in}%
\pgfsys@useobject{currentmarker}{}%
\end{pgfscope}%
\end{pgfscope}%
\begin{pgfscope}%
\definecolor{textcolor}{rgb}{0.000000,0.000000,0.000000}%
\pgfsetstrokecolor{textcolor}%
\pgfsetfillcolor{textcolor}%
\pgftext[x=2.616558in,y=0.365432in,,top]{\color{textcolor}{\rmfamily\fontsize{8.000000}{9.600000}\selectfont\catcode`\^=\active\def^{\ifmmode\sp\else\^{}\fi}\catcode`\%=\active\def%{\%}$\mathdefault{20}$}}%
\end{pgfscope}%
\begin{pgfscope}%
\pgfpathrectangle{\pgfqpoint{0.514740in}{0.462654in}}{\pgfqpoint{2.720000in}{0.807339in}}%
\pgfusepath{clip}%
\pgfsetbuttcap%
\pgfsetroundjoin%
\pgfsetlinewidth{0.501875pt}%
\definecolor{currentstroke}{rgb}{0.690196,0.690196,0.690196}%
\pgfsetstrokecolor{currentstroke}%
\pgfsetstrokeopacity{0.250000}%
\pgfsetdash{{1.850000pt}{0.800000pt}}{0.000000pt}%
\pgfpathmoveto{\pgfqpoint{3.111103in}{0.462654in}}%
\pgfpathlineto{\pgfqpoint{3.111103in}{1.269994in}}%
\pgfusepath{stroke}%
\end{pgfscope}%
\begin{pgfscope}%
\pgfsetbuttcap%
\pgfsetroundjoin%
\definecolor{currentfill}{rgb}{0.000000,0.000000,0.000000}%
\pgfsetfillcolor{currentfill}%
\pgfsetlinewidth{0.803000pt}%
\definecolor{currentstroke}{rgb}{0.000000,0.000000,0.000000}%
\pgfsetstrokecolor{currentstroke}%
\pgfsetdash{}{0pt}%
\pgfsys@defobject{currentmarker}{\pgfqpoint{0.000000in}{-0.048611in}}{\pgfqpoint{0.000000in}{0.000000in}}{%
\pgfpathmoveto{\pgfqpoint{0.000000in}{0.000000in}}%
\pgfpathlineto{\pgfqpoint{0.000000in}{-0.048611in}}%
\pgfusepath{stroke,fill}%
}%
\begin{pgfscope}%
\pgfsys@transformshift{3.111103in}{0.462654in}%
\pgfsys@useobject{currentmarker}{}%
\end{pgfscope}%
\end{pgfscope}%
\begin{pgfscope}%
\definecolor{textcolor}{rgb}{0.000000,0.000000,0.000000}%
\pgfsetstrokecolor{textcolor}%
\pgfsetfillcolor{textcolor}%
\pgftext[x=3.111103in,y=0.365432in,,top]{\color{textcolor}{\rmfamily\fontsize{8.000000}{9.600000}\selectfont\catcode`\^=\active\def^{\ifmmode\sp\else\^{}\fi}\catcode`\%=\active\def%{\%}$\mathdefault{25}$}}%
\end{pgfscope}%
\begin{pgfscope}%
\definecolor{textcolor}{rgb}{0.000000,0.000000,0.000000}%
\pgfsetstrokecolor{textcolor}%
\pgfsetfillcolor{textcolor}%
\pgftext[x=1.874740in,y=0.211111in,,top]{\color{textcolor}{\rmfamily\fontsize{8.000000}{9.600000}\selectfont\catcode`\^=\active\def^{\ifmmode\sp\else\^{}\fi}\catcode`\%=\active\def%{\%}Time ($\mu$s)}}%
\end{pgfscope}%
\begin{pgfscope}%
\pgfpathrectangle{\pgfqpoint{0.514740in}{0.462654in}}{\pgfqpoint{2.720000in}{0.807339in}}%
\pgfusepath{clip}%
\pgfsetbuttcap%
\pgfsetroundjoin%
\pgfsetlinewidth{0.501875pt}%
\definecolor{currentstroke}{rgb}{0.690196,0.690196,0.690196}%
\pgfsetstrokecolor{currentstroke}%
\pgfsetstrokeopacity{0.250000}%
\pgfsetdash{{1.850000pt}{0.800000pt}}{0.000000pt}%
\pgfpathmoveto{\pgfqpoint{0.514740in}{0.499527in}}%
\pgfpathlineto{\pgfqpoint{3.234740in}{0.499527in}}%
\pgfusepath{stroke}%
\end{pgfscope}%
\begin{pgfscope}%
\pgfsetbuttcap%
\pgfsetroundjoin%
\definecolor{currentfill}{rgb}{0.000000,0.000000,0.000000}%
\pgfsetfillcolor{currentfill}%
\pgfsetlinewidth{0.803000pt}%
\definecolor{currentstroke}{rgb}{0.000000,0.000000,0.000000}%
\pgfsetstrokecolor{currentstroke}%
\pgfsetdash{}{0pt}%
\pgfsys@defobject{currentmarker}{\pgfqpoint{-0.048611in}{0.000000in}}{\pgfqpoint{-0.000000in}{0.000000in}}{%
\pgfpathmoveto{\pgfqpoint{-0.000000in}{0.000000in}}%
\pgfpathlineto{\pgfqpoint{-0.048611in}{0.000000in}}%
\pgfusepath{stroke,fill}%
}%
\begin{pgfscope}%
\pgfsys@transformshift{0.514740in}{0.499527in}%
\pgfsys@useobject{currentmarker}{}%
\end{pgfscope}%
\end{pgfscope}%
\begin{pgfscope}%
\definecolor{textcolor}{rgb}{0.000000,0.000000,0.000000}%
\pgfsetstrokecolor{textcolor}%
\pgfsetfillcolor{textcolor}%
\pgftext[x=0.358489in, y=0.460947in, left, base]{\color{textcolor}{\rmfamily\fontsize{8.000000}{9.600000}\selectfont\catcode`\^=\active\def^{\ifmmode\sp\else\^{}\fi}\catcode`\%=\active\def%{\%}$\mathdefault{0}$}}%
\end{pgfscope}%
\begin{pgfscope}%
\pgfpathrectangle{\pgfqpoint{0.514740in}{0.462654in}}{\pgfqpoint{2.720000in}{0.807339in}}%
\pgfusepath{clip}%
\pgfsetbuttcap%
\pgfsetroundjoin%
\pgfsetlinewidth{0.501875pt}%
\definecolor{currentstroke}{rgb}{0.690196,0.690196,0.690196}%
\pgfsetstrokecolor{currentstroke}%
\pgfsetstrokeopacity{0.250000}%
\pgfsetdash{{1.850000pt}{0.800000pt}}{0.000000pt}%
\pgfpathmoveto{\pgfqpoint{0.514740in}{0.850538in}}%
\pgfpathlineto{\pgfqpoint{3.234740in}{0.850538in}}%
\pgfusepath{stroke}%
\end{pgfscope}%
\begin{pgfscope}%
\pgfsetbuttcap%
\pgfsetroundjoin%
\definecolor{currentfill}{rgb}{0.000000,0.000000,0.000000}%
\pgfsetfillcolor{currentfill}%
\pgfsetlinewidth{0.803000pt}%
\definecolor{currentstroke}{rgb}{0.000000,0.000000,0.000000}%
\pgfsetstrokecolor{currentstroke}%
\pgfsetdash{}{0pt}%
\pgfsys@defobject{currentmarker}{\pgfqpoint{-0.048611in}{0.000000in}}{\pgfqpoint{-0.000000in}{0.000000in}}{%
\pgfpathmoveto{\pgfqpoint{-0.000000in}{0.000000in}}%
\pgfpathlineto{\pgfqpoint{-0.048611in}{0.000000in}}%
\pgfusepath{stroke,fill}%
}%
\begin{pgfscope}%
\pgfsys@transformshift{0.514740in}{0.850538in}%
\pgfsys@useobject{currentmarker}{}%
\end{pgfscope}%
\end{pgfscope}%
\begin{pgfscope}%
\definecolor{textcolor}{rgb}{0.000000,0.000000,0.000000}%
\pgfsetstrokecolor{textcolor}%
\pgfsetfillcolor{textcolor}%
\pgftext[x=0.358489in, y=0.811958in, left, base]{\color{textcolor}{\rmfamily\fontsize{8.000000}{9.600000}\selectfont\catcode`\^=\active\def^{\ifmmode\sp\else\^{}\fi}\catcode`\%=\active\def%{\%}$\mathdefault{2}$}}%
\end{pgfscope}%
\begin{pgfscope}%
\pgfpathrectangle{\pgfqpoint{0.514740in}{0.462654in}}{\pgfqpoint{2.720000in}{0.807339in}}%
\pgfusepath{clip}%
\pgfsetbuttcap%
\pgfsetroundjoin%
\pgfsetlinewidth{0.501875pt}%
\definecolor{currentstroke}{rgb}{0.690196,0.690196,0.690196}%
\pgfsetstrokecolor{currentstroke}%
\pgfsetstrokeopacity{0.250000}%
\pgfsetdash{{1.850000pt}{0.800000pt}}{0.000000pt}%
\pgfpathmoveto{\pgfqpoint{0.514740in}{1.201549in}}%
\pgfpathlineto{\pgfqpoint{3.234740in}{1.201549in}}%
\pgfusepath{stroke}%
\end{pgfscope}%
\begin{pgfscope}%
\pgfsetbuttcap%
\pgfsetroundjoin%
\definecolor{currentfill}{rgb}{0.000000,0.000000,0.000000}%
\pgfsetfillcolor{currentfill}%
\pgfsetlinewidth{0.803000pt}%
\definecolor{currentstroke}{rgb}{0.000000,0.000000,0.000000}%
\pgfsetstrokecolor{currentstroke}%
\pgfsetdash{}{0pt}%
\pgfsys@defobject{currentmarker}{\pgfqpoint{-0.048611in}{0.000000in}}{\pgfqpoint{-0.000000in}{0.000000in}}{%
\pgfpathmoveto{\pgfqpoint{-0.000000in}{0.000000in}}%
\pgfpathlineto{\pgfqpoint{-0.048611in}{0.000000in}}%
\pgfusepath{stroke,fill}%
}%
\begin{pgfscope}%
\pgfsys@transformshift{0.514740in}{1.201549in}%
\pgfsys@useobject{currentmarker}{}%
\end{pgfscope}%
\end{pgfscope}%
\begin{pgfscope}%
\definecolor{textcolor}{rgb}{0.000000,0.000000,0.000000}%
\pgfsetstrokecolor{textcolor}%
\pgfsetfillcolor{textcolor}%
\pgftext[x=0.358489in, y=1.162969in, left, base]{\color{textcolor}{\rmfamily\fontsize{8.000000}{9.600000}\selectfont\catcode`\^=\active\def^{\ifmmode\sp\else\^{}\fi}\catcode`\%=\active\def%{\%}$\mathdefault{4}$}}%
\end{pgfscope}%
\begin{pgfscope}%
\definecolor{textcolor}{rgb}{0.000000,0.000000,0.000000}%
\pgfsetstrokecolor{textcolor}%
\pgfsetfillcolor{textcolor}%
\pgftext[x=0.302933in,y=0.866324in,,bottom,rotate=90.000000]{\color{textcolor}{\rmfamily\fontsize{8.000000}{9.600000}\selectfont\catcode`\^=\active\def^{\ifmmode\sp\else\^{}\fi}\catcode`\%=\active\def%{\%}$i_L$ (A)}}%
\end{pgfscope}%
\begin{pgfscope}%
\pgfpathrectangle{\pgfqpoint{0.514740in}{0.462654in}}{\pgfqpoint{2.720000in}{0.807339in}}%
\pgfusepath{clip}%
\pgfsetrectcap%
\pgfsetroundjoin%
\pgfsetlinewidth{1.505625pt}%
\definecolor{currentstroke}{rgb}{0.835294,0.368627,0.000000}%
\pgfsetstrokecolor{currentstroke}%
\pgfsetdash{}{0pt}%
\pgfpathmoveto{\pgfqpoint{0.638376in}{0.499525in}}%
\pgfpathlineto{\pgfqpoint{0.638416in}{0.499441in}}%
\pgfpathlineto{\pgfqpoint{0.638594in}{0.500164in}}%
\pgfpathlineto{\pgfqpoint{0.686919in}{0.709127in}}%
\pgfpathlineto{\pgfqpoint{0.803188in}{1.214253in}}%
\pgfpathlineto{\pgfqpoint{0.803207in}{1.214272in}}%
\pgfpathlineto{\pgfqpoint{0.803410in}{1.213815in}}%
\pgfpathlineto{\pgfqpoint{0.837989in}{1.133998in}}%
\pgfpathlineto{\pgfqpoint{0.882498in}{1.026983in}}%
\pgfpathlineto{\pgfqpoint{1.050643in}{0.617570in}}%
\pgfpathlineto{\pgfqpoint{1.085262in}{0.539190in}}%
\pgfpathlineto{\pgfqpoint{1.100098in}{0.506896in}}%
\pgfpathlineto{\pgfqpoint{1.105043in}{0.499352in}}%
\pgfpathlineto{\pgfqpoint{1.133097in}{0.499952in}}%
\pgfpathlineto{\pgfqpoint{1.182798in}{0.721482in}}%
\pgfpathlineto{\pgfqpoint{1.281708in}{1.163624in}}%
\pgfpathlineto{\pgfqpoint{1.297753in}{1.233297in}}%
\pgfpathlineto{\pgfqpoint{1.298114in}{1.232516in}}%
\pgfpathlineto{\pgfqpoint{1.327589in}{1.167193in}}%
\pgfpathlineto{\pgfqpoint{1.367153in}{1.075540in}}%
\pgfpathlineto{\pgfqpoint{1.421553in}{0.944906in}}%
\pgfpathlineto{\pgfqpoint{1.525407in}{0.694439in}}%
\pgfpathlineto{\pgfqpoint{1.564971in}{0.603045in}}%
\pgfpathlineto{\pgfqpoint{1.599589in}{0.526610in}}%
\pgfpathlineto{\pgfqpoint{1.609480in}{0.505582in}}%
\pgfpathlineto{\pgfqpoint{1.614425in}{0.499352in}}%
\pgfpathlineto{\pgfqpoint{1.627661in}{0.499996in}}%
\pgfpathlineto{\pgfqpoint{1.678309in}{0.725262in}}%
\pgfpathlineto{\pgfqpoint{1.777218in}{1.166425in}}%
\pgfpathlineto{\pgfqpoint{1.792299in}{1.231768in}}%
\pgfpathlineto{\pgfqpoint{1.792660in}{1.230985in}}%
\pgfpathlineto{\pgfqpoint{1.822135in}{1.165467in}}%
\pgfpathlineto{\pgfqpoint{1.861698in}{1.073589in}}%
\pgfpathlineto{\pgfqpoint{1.916098in}{0.942717in}}%
\pgfpathlineto{\pgfqpoint{2.019953in}{0.692023in}}%
\pgfpathlineto{\pgfqpoint{2.059516in}{0.600619in}}%
\pgfpathlineto{\pgfqpoint{2.094135in}{0.524216in}}%
\pgfpathlineto{\pgfqpoint{2.104026in}{0.503221in}}%
\pgfpathlineto{\pgfqpoint{2.108971in}{0.499352in}}%
\pgfpathlineto{\pgfqpoint{2.122207in}{0.499996in}}%
\pgfpathlineto{\pgfqpoint{2.172853in}{0.725306in}}%
\pgfpathlineto{\pgfqpoint{2.271762in}{1.166556in}}%
\pgfpathlineto{\pgfqpoint{2.286848in}{1.231966in}}%
\pgfpathlineto{\pgfqpoint{2.287135in}{1.231342in}}%
\pgfpathlineto{\pgfqpoint{2.319467in}{1.159355in}}%
\pgfpathlineto{\pgfqpoint{2.359030in}{1.067235in}}%
\pgfpathlineto{\pgfqpoint{2.413430in}{0.936209in}}%
\pgfpathlineto{\pgfqpoint{2.512339in}{0.697394in}}%
\pgfpathlineto{\pgfqpoint{2.556849in}{0.594611in}}%
\pgfpathlineto{\pgfqpoint{2.591467in}{0.518525in}}%
\pgfpathlineto{\pgfqpoint{2.596412in}{0.508027in}}%
\pgfpathlineto{\pgfqpoint{2.601358in}{0.499352in}}%
\pgfpathlineto{\pgfqpoint{2.616734in}{0.499952in}}%
\pgfpathlineto{\pgfqpoint{2.666441in}{0.721035in}}%
\pgfpathlineto{\pgfqpoint{2.770296in}{1.183943in}}%
\pgfpathlineto{\pgfqpoint{2.781393in}{1.231973in}}%
\pgfpathlineto{\pgfqpoint{2.781681in}{1.231350in}}%
\pgfpathlineto{\pgfqpoint{2.814012in}{1.159357in}}%
\pgfpathlineto{\pgfqpoint{2.853576in}{1.067230in}}%
\pgfpathlineto{\pgfqpoint{2.907976in}{0.936197in}}%
\pgfpathlineto{\pgfqpoint{3.006885in}{0.697370in}}%
\pgfpathlineto{\pgfqpoint{3.051394in}{0.594583in}}%
\pgfpathlineto{\pgfqpoint{3.086012in}{0.518494in}}%
\pgfpathlineto{\pgfqpoint{3.090958in}{0.507996in}}%
\pgfpathlineto{\pgfqpoint{3.095903in}{0.499352in}}%
\pgfpathlineto{\pgfqpoint{3.111103in}{0.499352in}}%
\pgfpathlineto{\pgfqpoint{3.111103in}{0.499352in}}%
\pgfusepath{stroke}%
\end{pgfscope}%
\begin{pgfscope}%
\pgfsetrectcap%
\pgfsetmiterjoin%
\pgfsetlinewidth{0.803000pt}%
\definecolor{currentstroke}{rgb}{0.000000,0.000000,0.000000}%
\pgfsetstrokecolor{currentstroke}%
\pgfsetdash{}{0pt}%
\pgfpathmoveto{\pgfqpoint{0.514740in}{0.462654in}}%
\pgfpathlineto{\pgfqpoint{0.514740in}{1.269994in}}%
\pgfusepath{stroke}%
\end{pgfscope}%
\begin{pgfscope}%
\pgfsetrectcap%
\pgfsetmiterjoin%
\pgfsetlinewidth{0.803000pt}%
\definecolor{currentstroke}{rgb}{0.000000,0.000000,0.000000}%
\pgfsetstrokecolor{currentstroke}%
\pgfsetdash{}{0pt}%
\pgfpathmoveto{\pgfqpoint{3.234740in}{0.462654in}}%
\pgfpathlineto{\pgfqpoint{3.234740in}{1.269994in}}%
\pgfusepath{stroke}%
\end{pgfscope}%
\begin{pgfscope}%
\pgfsetrectcap%
\pgfsetmiterjoin%
\pgfsetlinewidth{0.803000pt}%
\definecolor{currentstroke}{rgb}{0.000000,0.000000,0.000000}%
\pgfsetstrokecolor{currentstroke}%
\pgfsetdash{}{0pt}%
\pgfpathmoveto{\pgfqpoint{0.514740in}{0.462654in}}%
\pgfpathlineto{\pgfqpoint{3.234740in}{0.462654in}}%
\pgfusepath{stroke}%
\end{pgfscope}%
\begin{pgfscope}%
\pgfsetrectcap%
\pgfsetmiterjoin%
\pgfsetlinewidth{0.803000pt}%
\definecolor{currentstroke}{rgb}{0.000000,0.000000,0.000000}%
\pgfsetstrokecolor{currentstroke}%
\pgfsetdash{}{0pt}%
\pgfpathmoveto{\pgfqpoint{0.514740in}{1.269994in}}%
\pgfpathlineto{\pgfqpoint{3.234740in}{1.269994in}}%
\pgfusepath{stroke}%
\end{pgfscope}%
\begin{pgfscope}%
\definecolor{textcolor}{rgb}{0.000000,0.000000,0.000000}%
\pgfsetstrokecolor{textcolor}%
\pgfsetfillcolor{textcolor}%
\pgftext[x=0.569140in,y=1.213480in,left,top]{\color{textcolor}{\rmfamily\fontsize{8.000000}{9.600000}\selectfont\catcode`\^=\active\def^{\ifmmode\sp\else\^{}\fi}\catcode`\%=\active\def%{\%}(b)}}%
\end{pgfscope}%
\end{pgfpicture}%
\makeatother%
\endgroup%

    \caption{(a) tensão de saída e (b) corrente do indutor do circuito Buck.}
    \label{fig:ex1}
\end{figure}